\section{Numerical Results and Discussions}
\label{sec:numerical}

\begin{figure*}[htb]
    \centering
    \includegraphics[width=0.48\textwidth]{images/O0IB_ncls_W3_Hybrid_ME_lat_pheno.pdf}
    \includegraphics[width=0.48\textwidth]{images/O0IB_ncls_H5_Hybrid_ME_lat_pheno.pdf}\\
    \includegraphics[width=0.48\textwidth]{images/O0IB_ncll_W3_Hybrid_ME_lat_pheno.pdf}
    \includegraphics[width=0.48\textwidth]{images/O0IB_ncll_H5_Hybrid_ME_lat_pheno.pdf}
    \caption{\UPDATE{Lattice (points) hybrid renormalized nucleon matrix elements with $z_s = 0.36$~fm for the Wilson-3 (left column) and HYP5 (right column) smearing for $M_\pi \approx 690$ (top row) and $310$ (bottom row) GeV compared to the matrix elements reconstructed from the CT18 gluon PDF~\cite{Hou:2019efy}, shown by the colored bands which have the same color scheme as the lattice data for differing momenta.}}
    \label{fig:hybrid-lat-pheno}
\end{figure*}

We use the high-statistics data from Ref.~\cite{Good:2024iur} on one ensemble with lattice spacing $a \approx 0.12$~fm at valence pion masses $M_{\pi} \approx 310$ and $690$~MeV generated by the MILC collaboration~\cite{MILC:2013znn} using 2+1+1 flavors of highly improved staggered quarks (HISQ)~\cite{Follana:2007rc} with the lattice volume of $24^3 \times 64$.
We use Wilson clover fermions in the valence sector and tune the valence quark mass to reproduce the light and strange mass in the HISQ sea.
We computed 1,296,640 two-point correlators across 1013 configurations for both the nucleon and pion at the light and heavy pion masses with fixed Gaussian momentum smearing~\cite{Bali:2016lva}.
The nucleons and pions are referred to as the light ($N_l$) and strange ($N_s$) nucleons, the pion ($\pi$), and $\eta$-strange ($\eta_s$) for the light and heavy pion masses, respectively.
We will focus mostly on the nucleons but we include the pion data to understand how the external states affect the calculation of $\delta m + m_0$.
We measured two sets of gluon loops with 5 steps of hypercubic \UPDATE{\cite{Hasenfratz:2001hp}} (HYP5) smearing and with Wilson flow\UPDATE{\cite{Luscher:2010iy}} with a flow time of $T=3 a^2$ (Wilson3) on the gauge links \UPDATE{in order to achieve reasonable signal.
A more detailed study of smearing on the matrix element level is provided by Ref.~\cite{NieMiera:2025inn} for similar kaon gluon data.
A key comparison here is that the ratio of renormalized Wilson3 smearing is similar to about 10 steps of HYP smearing.}

With the zero momentum matrix elements, $h^\text{B}(z,0)$, and the Wilson coefficient defined in Eq.~\ref{eq:wil-coef}, we may fit $\delta m + m_0$ using the form in Eq.~\ref{eq:dm_fit}.
The $a\approx 0.12$~fm lattice spacing is too coarse to capture where $\ln{\left[\mathcal{H}(z,\mu) /h^{\text{B}}(z,0) \right]}$ behaves linearly in the small-$z$ region, so we interpolate the $h^\text{B}(z,0)$.
We fit the interpolated data in the set $\{z-0.2, z, z+0.2\}$ in units of fm, varying $z$.
\UPDATE{We find that the results are highly consistent for any order of interpolation between 3 and 8.}
We show these results for the $\delta m + m_0$ versus $z$ for each hadron and each smearing, in Fig.~\ref{fig:m0_vs_z} with statistical errors only.
We do not consider scale variation here.
The general behavior of $\delta m + m_0$ is very reasonable here for all cases.
The fits are unstable in the small-$z$ region due to discretization effects, then there is a nice plateau of stability, before we see the hint of divergence due to perturbation theory breaking down as $z$ becomes larger.
As seen by comparing the left and right plots, the amount of smearing can have a significant effect on the value of $\delta m + m_0$.
Focusing on the higher precision Wilson3 data on the left plot, we see that $\delta m + m_0$ has very little dependence on the external state, with hadrons at the same pion mass being statistically indistinguishable and the same hadrons at different pion masses having agreement well within about $2\sigma$.
These results are in contrast to those from the other two operators explored in Fig. 11 of Ref.~\cite{Good:2024iur}.
For those operators, the fitted $\delta m + m_0$ values depend highly on the external states and do not appear to level off as nicely.
This adds even further to the list of reasons why the operator in this work is preferable to the other two.
We choose $\text{min}_z(\delta m + m_0)$ as the value to use for hybrid renormalization moving forward for each hadron and smearing, as these happen to be where the curves are the flattest.
We note that, for Wilson-3, these values are between $0.35$ and $0.4$ GeV, which is slightly smaller than the value $0.5$ GeV, which we estimated in our previous paper without knowing the Wilson coefficients.

\begin{figure}[t]
    \centering
    \includegraphics[width=\linewidth]{images/O0IBncls_W3_large-nu.pdf}
    \caption{Large-$\nu$ extrapolation of the hybrid renormalized $P_z = 1.71$ GeV matrix elements \UPDATE{with $M_\pi \approx 690$ MeV and Wilson3 smearing}, using the data from $z \in [7a, 16a]$, the extent of which is shown by the two dashed, vertical lines.}
    \label{fig:O0IB_large-nu}
\end{figure}

We plot the hybrid renormalized matrix elements using the fitted $\delta m + m_0$ values from the zero momentum matrix elements for the light and strange nucleons for the Wilson-3 and HYP5 smearing along with the nucleon matrix elements reconstructed from the CT18 NNLO gluon PDF~\cite{Hou:2019efy} in Fig.~\ref{fig:hybrid-lat-pheno} \UPDATE{to provide a qualitative example of matrix elements from a physically reasonable PDF.}
\UPDATE{These matrix elements are computed by using the CT18 gluon PDF on the righthand side of Eq.~\ref{eq:mtcheq} to obtain the quasi-PDF, then Fourier transforming back to position space.}
We use a value of $z_s = 3a \approx 0.36$~fm.
Looking at the lattice data first, as one would expect from increased gauged link smearing, we see that the Wilson-3 matrix elements seem to fall off more slowly than the HYP5 matrix elements.
This suggests that the Wilson-3 smearing is affecting the physics, especially when we compare with the phenomenological matrix elements.
We see that the low momentum matrix elements seem to fall off much faster and disagree more significantly with the phenomenological matrix elements, with the largest momentum matrix elements beginning to converge more quickly in the strange nucleon case.
We are not overly concerned with the lowest momenta data because we intend to apply LaMET matching only to large-momentum data, but this may be interesting to explore further to better understand the systematics.
It seems that the HYP5 strange nucleon is in the best agreement with the phenomenological results, while the ``most physical" data should be from the light nucleon for HYP5.
This could be some canceling of various systematic effects, which will all require further exploration.
One further curiosity is that there is still a bump near $\nu_s = z_sP_z$ in the phenomenological matrix elements for the lowest momentum that was seen in our previous results for the other two gluon operators in Fig. 14 of Ref.~\cite{Good:2024iur}; however, it is much smaller and not very present in the larger momentum results at all.
This is due to the differences in the matching for the different operators and has nothing to do with the lattice value of $\delta m + m_0$.

\begin{figure}
    \centering
    \includegraphics[width=\linewidth]{images/qPDF-PDF-comp_H5_ncll_O0IB_Pz1.71.pdf}
    \caption{Quasi-PDF (green) from the Fourier transform of the hybrid renormalized $P_z = 1.71$ GeV matrix elements with $M_\pi \approx 310$ MeV and HYP5 smearing and the light-front (LF-)PDF (purple) after matching the quasi-PDF (qPDF) to the light cone.
    The gray bands represent the region where LaMET matching begins to break down.}
    \label{fig:qPDF-PDF-comp}
\end{figure}

\begin{figure*}[ht!]
    \centering
    \includegraphics[width=.45\textwidth]{images/AllComp_Pz1.71.pdf}
    \includegraphics[width=.45\textwidth]{images/CT18comp_H5_O0IB_Pz1.71.pdf}
    \caption{ (Left) the MSULat nucleon gluon PDFs at $P_z = 1.71$ for the different pion masses and smearings all compared.
    (Right) The MSULat PDF  from large momentum effective theory at $P_z = 1.71$ GeV at $M_\pi \approx 690$ (blue) and $310$ (red) MeV with HYP5 smearing compared to the CT18~\cite{Hou:2019efy} (orange), JAM24~\cite{Anderson:2024evk} (cyan), and CJ22 with multiplicative higher twist corrections~\cite{Cerutti:2025yji} (green) gluon PDFs.
    In both plots, the gray bands represent the region where LaMET matching begins to break down.
    }
    \label{fig:PDF-GF-comp}
\end{figure*}

Thanks to the extremely high statistics, we have reasonable signal at large enough distances and high enough momentum to apply a large-$\nu$ extrapolation to do a Fourier transformation.
We follow previous work and use an ansatz of the form~\cite{Ji:2020brr}
%
\begin{equation}\label{eq:large-nu-form}
    h^\text{R}(z,P_z) \approx A\frac{e^{-m\nu}}{|\nu|^d}
\end{equation}
%
to fit the large-$\nu$ data, where $A$, $m$ and $d$ are fitted parameters.
We use this form to fit our nucleon data at $P_z = 1.71$~GeV data for each smearing and pion mass.
\UPDATE{This fit form has strong theoretical motivation described in the supplemental material of Ref.~\cite{Gao:2021dbh}, which allows us to replace the noisy long-distance data with the well constrained, physically motivated extrapolation.
This has been compared to standard lattice methods of using theoretically motivated fit forms to extract information from limited data, justifying the truncation of noisy data.
One such example is extracting matrix elements from lattice correlator data, which is described in our earlier paper~\cite{Good:2024iur}.}
The fit window for the extrapolation that we used in our previous work~\cite{Good:2024iur} was overly conservative, as we can fit data as low as $z = 6a = 0.72$ fm, and still achieve consistent results, which is explored further in Appendix~\ref{sec:app_large-nu-comp}.
\UPDATE{The results imply that there is systematic uncertainty associated with the fit range.
We do not explore this systematic in detail in this first proof of principle calculation.}
We use fit windows that start at $6a$ for the HYP5 data, and $7a$ for the Wilson3 data.
We find again that it is still difficult to separate the algebraic and exponential decay at this level of precision, which causes a very large amount of fluctuation in the fit parameters, but this fluctuation of the individual parameters cancels and allows us to obtain a reasonable fit and extrapolation.

We plot the results of this fit for the Wilson3 strange nucleon in Fig.~\ref{fig:O0IB_large-nu}, as an example of our most precise data.
We see qualitatively that the fit agrees well with the data and decays nicely to zero while the largest distance data diverge \UPDATE{due to the increasing UV fluctuations and finite volume effects introduced at such large distances}.
Inside the fit range, the parameterization has error similar to the data, before shrinking exponentially compared to the exponentially growing error in the data.
This is consistent for the other datasets, which are shown in Appendix~\ref{sec:app_large-nu}.
With these two features, we can reliably perform a Fourier transformation on the data.

We can connect an interpolation of the short distance lattice matrix elements and the large-$\nu$ extrapolation at longer distances to compute the Fourier transformation of the matrix elements, which defines the quasi-PDF.
We apply the light-front matching defined in Eqs.~\ref{eq:mtcheq} and \ref{eq:kernel} to obtain the light-front PDF.
As an example of the other end of the spectrum of our statistical precision, Fig.~\ref{fig:qPDF-PDF-comp} shows the results for the quasi- and light-front-PDFs for the light nucleon with HYP5 smearing, with the regions with $\frac{\Lambda_\text{QCD}}{P_zx(1-x)} \gtrsim 1$, where LaMET matching breaks down, grayed out.
\UPDATE{We use $\Lambda_\text{QCD} = 300$ MeV as a conservative estimate.}
Focusing first on the quasi-PDF in green, we see that there is some negativity in the quasi-PDF which was also seen in our previous estimation in Ref.~\cite{Good:2024iur}.
We also see pinch points in the errors on the quasi-PDF, which is likely due to the fairly large range in the decay rates found in our large-$\nu$ extrapolation.
This is most significant for the HYP5 light nucleon, which we show here.
A nice feature is that the quasi-PDF is very well convergent to 0 as $x \rightarrow 1$.
Comparing the light-front- and quasi-PDFs, we can see that the matching has a significant effect on the PDFs.
This seems to be larger than what is often seen in the quark PDFs, but this can be partially accounted for due to the difference in the factors $C_F = 4/3$ and $C_A = 3$ in front of the NLO term in the matching for the quark and gluon PDFs, respectively.
The negative portion of the quasi-PDF is mostly removed by the matching, with the light-front PDF having a slightly negative central value from $x \gtrsim 0.75$, which is mostly within $1\sigma$ of zero.
Another striking effect of the matching is that the errors seem to be significantly reduced.
It could be that the oscillation in the errors of the quasi-PDF undergoes some level of destructive interference through the matching or that the matching is simply very robust to the statistical fluctuations in the data.
It will be interesting to explore this further as the data becomes more precise.

To get a picture of the systematics from smearing and pion masses, we show all of our nucleon gluon PDFs from the $P_z = 1.71$~GeV lattice data at the two pion masses and two smearings with only statistical fluctuations on the left side of Fig.~\ref{fig:PDF-GF-comp}.
The Wilson3 smearing data produces PDFs that are more negative at $x \gtrsim 0.6$, but within about 1 or 2$\sigma$ of zero.
At $x \approx 0.5$, where LaMET theory is most reliable, the $M_\pi \approx 690$ MeV PDFs are around $3\sigma$ from each other, which demonstrates that there are significant effects on the physics caused by the very large Wilson3 smearing.
\UPDATE{In particular, this effect in the large-$x$ may be coming from the effects of smearing on the short-distance physics or perhaps imperfect renormalization, which has not been proven analytically for smeared lattice studies.}
Overall, we believe that the HYP5 PDFs are more reliable than the Wilson3 PDFs.
The PDFs at the two different pion masses for each smearing are within one sigma of each other, with those at the smaller pion mass having larger statistical fluctuations, as one would expect.
There is still more work to be done in understanding how the systematic effects emerge in the final PDFs, but this is not the focus of this first study.
\UPDATE{It would be interesting to explore renormalizability in the continuum theory with Wilson flow.}

We compare our HYP5 PDFs at each pion mass to the CT18~\cite{Hou:2019efy}, JAM24~\cite{Anderson:2024evk}, and the CJ22 multiplicative higher twist effect~\cite{Cerutti:2025yji} gluon PDFs, which are fairly different global fits from each other, in Fig.~\ref{fig:PDF-GF-comp}.
We see that the HYP5 lattice PDFs actually agree very well with the CT18 across the entire region, with the strange nucleon agreeing particularly well, which is expected from the agreement seen in position space in Fig.~\ref{fig:hybrid-lat-pheno}.
We see that on the lower and upper side of the region that we trust, we see 1$\sigma$ agreement with the CJ22 PDFs, while it seems that the agreement is lesser around $x\approx 0.5$ due to the pinch point in the lattice error bars here.
We see some overlap with the JAM24 PDF at several different points in the trustworthy region, with particular agreement with the lighter pion mass, partially due to the larger errors.
Overall, our PDF falls within the range of global fits, preferring a slightly larger gluon PDF around $x\approx 0.5$ than most global fits.
\UPDATE{This provides a loose validation that the methodology can provide qualitatively reasonable gluon PDF; however, there are many systematic effects to be understood.
At this level, it is not abundantly clear whether there is some cancellation of systematic effects causing the agreement.}
It will be interesting to see what occurs as noise reduction techniques reduce the error on the lighter pion mass results and other lattice systematics are taken into consideration, such as physical continuum extrapolations.
Finer lattice spacings will be needed to go to higher momenta, as well as to increase the reliable region of our PDFs.
